\section{Lab 11: Events and Delegates in C\#}

\subsection*{Metadata}
\begin{itemize}
	\item Course: CS202 - Software Tools and Techniques for CSE
	\item Lab Topic: Events and Delegates in C\# Windows Forms Applications
	\item Name: Shardul Junagade (23110297)
	\item Date: 3rd November 2025
\end{itemize}

Repository path: \texttt{lab11/}

\subsection{Introduction}
I completed Laboratory Session 11 on events and delegates in C\# Windows Forms. The session centered on building a small interactive form and then reasoning through several delegate/event code snippets to understand chaining, thresholds, and nested invocation. In this lab I focused on wiring custom events between UI controls, extending a basic event with a custom \texttt{EventArgs} payload, and demonstrating multicast behavior. After that I manually traced five embedded C\# code snippets (Tasks 3 - 5) and wrote out their outputs with reasoning. I kept the implementation compact and readable rather than over-engineered.

\subsection{Tools}
\begin{itemize}[itemsep=0pt]
	\item OS: Windows 11; Shell: PowerShell 7
	\item .NET SDK: .NET 8 (meets the ``.NET 6 or later'' requirement)
	\item IDE/Editor: Visual Studio 2022 Community
	\item Git: for version control
\end{itemize}

\insertimage{./images/env.png}[Environment Details]

\subsection{Setup}
I reused the same .NET setup as in earlier labs and verified it before working on the Windows Forms app.

\subsubsection*{Visual Studio path}
\begin{enumerate}
	\item Install Visual Studio 2022 (Community Edition) with the ``.NET desktop development'' workload selected.

	\insertimage{./images/workload.png}[Workload Selection]

	\item Confirm the SDK is on PATH:
\begin{verbatim}
dotnet --version
\end{verbatim}
	\item Create a new C\# Windows Forms App in Visual Studio and set the Target Framework to .NET 8.0 in project properties.

	\insertimage{./images/1.png}[Project Target Framework]

	\insertimage{./images/net8.png}[.NET Version Selection]
\end{enumerate}


\subsection{Methodology and Execution}

I then designed the form with two buttons (\texttt{btnChangeColor}, \texttt{btnChangeText}), one label (\texttt{lblMessage} initialised to ``Welcome to Events Lab''), and a ComboBox (\texttt{cmbColors}) where I added the entries Red, Green, Blue via the Properties window.

\insertimage{./images/2.0.png}[Form Design]

\subsubsection{Part 1 - Basic Custom Events (Color \& Text)}
In \texttt{Form1.cs} I declared two delegates and the corresponding events:
\begin{lstlisting}
public delegate void ColorChangedHandler(System.Drawing.Color newColor);
public delegate void TextChangedHandler(string newText);
public event ColorChangedHandler ColorChangedEvent;
public event TextChangedHandler TextChangedEvent;
\end{lstlisting}

In the form constructor I subscribed the handlers:
\begin{lstlisting}
ColorChangedEvent += UpdateLabelColor;
TextChangedEvent += UpdateLabelText;
\end{lstlisting}

On the colour button click I read the selected ComboBox string, mapped it to a \texttt{System.Drawing.Color} with a \texttt{switch}, and invoked the custom event:
\begin{lstlisting}
ColorChangedEvent?.Invoke(selectedColor); // raises custom event
\end{lstlisting}

The text button built a timestamp string and invoked:
\begin{lstlisting}
TextChangedEvent?.Invoke(newText);
\end{lstlisting}

The label foreground and caption updated only through these events - the click handlers acted as publishers, and the subscriber methods performed the UI changes.

\insertimage{./images/2.1.png}[Part 1 - Code Snippet of Form1.cs]

\insertimage{./images/2.2.png}[Part 1 - Code Snippet of Form1.cs]

\insertimage{./images/3.png}[Part 1 - GUI after colour change]

\insertimage{./images/4.png}[Part 1 - GUI after text change]

\subsubsection{Part 2 - EventArgs + Multicast}
For the second part I extended the colour change event to carry structured information. I added a \texttt{ColorEventArgs} class:
\begin{lstlisting}
public class ColorEventArgs : EventArgs
{
	public string ColorName { get; }

	public ColorEventArgs(string colorName)
	{
		ColorName = colorName;
	}
}
\end{lstlisting}

\insertimage{./images/5.png}[ColorEventArgs class]

Then in \texttt{Form1.cs} I adopted the conventional sender + args pattern for the event:
\begin{lstlisting}
public delegate void ColorChangedHandler(object sender, ColorEventArgs e);
public event ColorChangedHandler ColorChangedEvent;
\end{lstlisting}

In the constructor I wired two subscribers to demonstrate multicast behaviour:
\begin{lstlisting}
ColorChangedEvent += UpdateLabelColor;
ColorChangedEvent += ShowNotification; // second subscriber
\end{lstlisting}

On the button click I created the EventArgs and invoked the event:
\begin{lstlisting}
var args = new ColorEventArgs(selectedColorName);
ColorChangedEvent?.Invoke(this, args);
\end{lstlisting}

	\texttt{UpdateLabelColor} mapped the \texttt{ColorName} string to a \texttt{System.Drawing.Color}, and \texttt{ShowNotification} popped up a message box confirming the choice. Both ran in order when I clicked the button, this showed multiple subscribers on the same event.

\insertimage{./images/6.0.png}[ColorEventArgs code snippet]

\insertimage{./images/6.1.png}[Part 2 - Code Snippet of Form1.cs]

\insertimage{./images/6.2.png}[Part 2 - Code Snippet of Form1.cs]

\insertimage{./images/6.3.png}[Part 2 - GUI and message box after colour change]

\subsubsection{Task 3 - Output Reasoning (Level 0)}

\paragraph{Q1.} Final output: 
\begin{lstlisting}
MS:-1
\end{lstlisting}

\insertimage{./images/7.png}[Task 3 Q1 - code and output]

Reasoning: the delegate started with \texttt{Add}, then had \texttt{Mul} and \texttt{Sub} added, and finally \texttt{Add} was removed. That left the invocation list as \texttt{[Mul, Sub]}. Calling \texttt{Mul(2,3)} printed \texttt{M} and returned 6; calling \texttt{Sub(2,3)} printed \texttt{S} and returned -1. For multicast delegates with a return type, only the value from the last method is used, so the printed result became \texttt{:-1}.

\newpage

\paragraph{Q2.} Final output: 
\begin{lstlisting}
I5 D4 F4
\end{lstlisting}

\insertimage{./images/8.png}[Task 3 Q2 - code and output]

Reasoning: with \texttt{val = 3}, the invocation list was \texttt{[Inc, Dec]}. \texttt{Inc(ref val)} increased it to 5 and printed \texttt{I5 }. Then \texttt{Dec(ref val)} dropped it to 4 and printed \texttt{D4 }. After the multicast call the variable held 4, so the final print produced \texttt{F4}.

\newpage

\subsubsection{Task 4 - Output Reasoning (Level 1)}

\paragraph{Q1.} Final output:
\begin{lstlisting}
>1>2[M2]>3[L3](Reset)>4[M4{Alert}]>5>6[M6][L6](Reset)
\end{lstlisting}

\insertimage{./images/9.png}[Task 4 Q1 - code and output]

Reasoning: the \texttt{Increment()} method first increases the internal counter and immediately prints the current value in the form \texttt{>value}. After that, two conditional checks determine whether additional text is printed.

If the value is even (\texttt{value \% 2 == 0}), the \texttt{MilestoneReached} event is raised, and its subscriber prints \texttt{[Mvalue]}. For the special case where the value is 4, the same subscriber also appends \texttt{\{Alert\}}. If the value is a multiple of 3 (\texttt{value \% 3 == 0}), the \texttt{LimitReached} event is raised, and two subscribers print \texttt{[Lvalue]} and \texttt{(Reset)} in that order.

The loop in \texttt{Main} calls \texttt{Increment()} six times, so the counter takes the values 1 through 6, and the printed output for each call is as follows:
\begin{itemize}[itemsep=0pt]
	\item Value 1: only \texttt{>1} (no event conditions satisfied).
	\item Value 2: \texttt{>2} and \texttt{[M2]} (even, milestone only).
	\item Value 3: \texttt{>3[L3](Reset)} (multiple of 3, limit reached only).
	\item Value 4: \texttt{>4[M4\{Alert\}]} (even, milestone with alert; not a multiple of 3).
	\item Value 5: only \texttt{>5} (no event conditions satisfied).
	\item Value 6: \texttt{>6[M6][L6](Reset)} (both even and multiple of 3, so milestone and limit events fire).
\end{itemize}

Concatenating these pieces in order produces the final output:
\begin{lstlisting}
>1>2[M2]>3[L3](Reset)>4[M4{Alert}]>5>6[M6][L6](Reset)
\end{lstlisting}

\newpage

\paragraph{Q2.} Final output:
\begin{lstlisting}
Temperature changed from 30 C to 46 C
 Warning: Sudden change detected!
Temperature changed from 46 C to 52 C
\end{lstlisting}

\insertimage{./images/10.png}[Task 4 Q2 - code and output]

Reasoning: the event fired only when the absolute temperature change was greater than 5. The changes 25→28 and 28→30 were ignored. The jump 30→46 (difference 16) triggered both the main message and the warning because the difference exceeded 10. The jump 46→52 (difference 6) triggered only the main message.

\newpage

\subsubsection{Task 5 - Output Reasoning (Level 2)}

\paragraph{Q1.} Final output:
\begin{lstlisting}
[Start]{Ping}(Nested)[Start]{Pong}(Nested)[End][End]
\end{lstlisting}

\insertimage{./images/11.png}[Task 5 Q1 - code and output]

Reasoning: calling \texttt{Trigger("Ping")} first printed \texttt{[Start]}. The first subscriber printed \texttt{\{Ping\}}, and the second printed \texttt{(Nested)} and then invoked \texttt{Trigger("Pong")} because the message was ``Ping''. The nested trigger printed another \texttt{[Start]}, then the two subscribers printed \texttt{\{Pong\}} and \texttt{(Nested)}. Since the message was now ``Pong'', there was no further recursion. Each \texttt{Trigger} call ended by printing \texttt{[End]}, so two \texttt{[End]} markers appeared at the end.

\newpage

\paragraph{Q2.} Final output: 
\begin{lstlisting}
[Check]{High}[Check][Done](Alert)[Done]
\end{lstlisting}

\insertimage{./images/12.png}[Task 5 Q2 - code and output]

Reasoning: \texttt{Check(80)} printed \texttt{[Check]} and then fired the threshold event because 80 was greater than 50. The first subscriber printed \texttt{\{High\}} and called \texttt{Check(30)} on the same sensor, which printed \texttt{[Check][Done]} but did not raise the event again because 30 $\leq$ 50. The second subscriber then printed \texttt{(Alert)}. After both subscribers finished, the original \texttt{Check(80)} call printed its final \texttt{[Done]}.

\subsection{Conclusion}
In this lab I built a Windows Forms example that used custom delegates and events to drive UI updates rather than direct control manipulation inside button handlers. I then extended the design using a custom \texttt{ColorEventArgs} and demonstrated multicast events by wiring both a label update and a message box to the same publisher. Overall the exercise strengthened my comfort with the publisher-subscriber pattern in C\# and clarified practical details like passing \texttt{ref} parameters through multicast delegates, and composing EventArgs for extensibility without changing event signatures.

\subsection{References}
\begin{itemize}[itemsep=0pt]
	\item Microsoft C\# docs: \url{https://learn.microsoft.com/en-us/dotnet/csharp}
	\item Delegates and events: \url{https://learn.microsoft.com/en-us/dotnet/csharp/programming-guide/delegates-and-events}
	\item Windows Forms events: \url{https://learn.microsoft.com/en-us/dotnet/desktop/winforms/}
\end{itemize}
